La transformation numérique constitue aujourd'hui un levier majeur d'amélioration de la performance des organisations, y compris dans les établissements publics. En automatisant les tâches répétitives, en fiabilisant les données et en offrant une vision consolidée de l'activité, les applications de gestion permettent de gagner en efficacité, en traçabilité et en qualité de service.

C'est dans ce contexte que s'inscrit le présent stage, réalisé au sein de l'Office National des Pêches (ONP). L'ONP accompagne le développement du secteur de la pêche au Maroc et met notamment à disposition des professionnels de la glace en blocs, indispensable à la conservation du poisson depuis le débarquement jusqu'à la commercialisation. La vente de cette glace génère une activité récurrente qui, dans le cadre de ce projet, était suivie de façon essentiellement manuelle.

Cette gestion manuelle présente plusieurs limites : risque d'erreurs de saisie et de calcul, absence de numérotation fiable des tickets de vente, difficulté à retrouver l'historique des transactions, et impossibilité de disposer d'indicateurs synthétiques sur l'activité. L'objectif du stage est donc de \textbf{concevoir et réaliser une application web} qui centralise et sécurise la gestion des ventes de blocs de glace, tout en restant simple à utiliser au quotidien par les agents et caissiers.

L'application développée, baptisée \textbf{IceVentes}, couvre la gestion des utilisateurs et de leurs rôles, le paramétrage du prix de la glace, la gestion des clients, la saisie des ventes avec génération automatique d'un code unique, ainsi que la consultation de l'historique et le pilotage de l'activité à travers un tableau de bord.

Ce rapport est organisé en quatre chapitres :

\begin{itemize}
  \item le \textbf{premier chapitre} présente le cadre du projet : l'organisme d'accueil, le contexte, la problématique, les objectifs et la conduite du projet ;
  \item le \textbf{deuxième chapitre} est consacré à la spécification des besoins fonctionnels et non fonctionnels, ainsi qu'à l'identification des acteurs et des cas d'utilisation ;
  \item le \textbf{troisième chapitre} détaille la conception du système : architecture, modèle de données et diagrammes de séquence ;
  \item le \textbf{quatrième chapitre} décrit la réalisation : environnement technique, interfaces produites, fonctionnalités clés et validation.
\end{itemize}

Le rapport se termine par une conclusion générale qui dresse le bilan du travail réalisé et propose des perspectives d'évolution.
